% sections/05_evaluation.tex

\section{Evaluation Framework}
\label{sec:evaluation}

\subsection{Layer 1: Deterministic Harness}

All quantitative results in this paper are produced by the Layer-1 evaluation
harness (\texttt{src/chemometrics\_mcp/core/evaluation.py}).
Layer 1 is defined by three properties:
\begin{itemize}
  \item \textbf{No LLM.} Tools are called directly as Python functions via
    \texttt{from chemometrics\_mcp.tools import \ldots}; no MCP transport,
    no language model, and no prompt is involved.
  \item \textbf{Fixed seed.} All stochastic operations (random splits,
    bootstrapping, shuffles) are seeded with \texttt{seed=42}.
  \item \textbf{Byte-reproducible.} Given the same seed and the same input
    files, every tool call returns bit-identical JSON output.
    This is verified by running \texttt{run\_scenario\_full} twice and
    asserting equality of the serialized output dicts.
\end{itemize}

Layer~2 evaluation---measuring system behavior with a live LLM issuing tool
calls through the MCP transport---is explicitly deferred to future work.
No Layer-2 result numbers appear in this paper.

\subsection{Evaluation Scenarios}

Table~\ref{tab:scenarios} describes the four evaluation scenarios.

\begin{table}[ht]
  \centering
  \caption{Evaluation scenario matrix. $n$: total spectra after loading.
    $p$: spectral features. Task: regression (Reg) or classification (Cls).
    Groups: number of distinct group labels.}
  \label{tab:scenarios}
  \begin{tabular}{llllll}
    \toprule
    ID & Name & Modality & Task & $n$ & Groups \\
    \midrule
    S1 & Wear-layer regression   & NIR  & Reg  & TODO & TODO \\
    S2 & Species multiclass      & NIR  & Cls  & TODO & TODO \\
    S3 & Vinyl-vs-wood binary    & NIR  & Cls  & TODO & TODO \\
    S4 & FTIR purity multiclass  & FTIR & Cls  & $\sim$20 & 7 \\
    \bottomrule
  \end{tabular}
\end{table}

\paragraph{S1 -- NIR wear-layer regression.}
Predict vinyl flooring wear-layer thickness (mil) from NIR spectra
collected in the range 1450--2450~nm.
The workbook contains replicates keyed by sample filename stem.
The task is \texttt{regression}; primary metric is RMSE and $R^2$.

\paragraph{S2 -- NIR species multiclass.}
Classify lumber species from the same NIR workbook.
Multiple replicate scans exist per species instance.
The task is \texttt{multi\_class\_classification}; primary metric is accuracy.

\paragraph{S3 -- NIR vinyl-vs-wood binary.}
Binary classification distinguishing vinyl from wood/lumber samples using
the same NIR workbook.
This is the simplest classification task (two well-separated classes
in NIR space) and serves as a sanity check that the pipeline produces
reasonable outputs.

\paragraph{S4 -- FTIR purity multiclass.}
Classify compound purity groups from FTIR transmittance spectra
($\approx$20 spectra from 7 measurement groups, collected June 2026).
This is the hardest scenario: small $n$, structured groups, and real
instrument variation across groups.

\subsection{Evaluation Metrics}

The harness measures seven metric categories:

\begin{enumerate}
  \item \textbf{HITL compliance.}
    Fraction of tool calls in which \texttt{approval\_required=True} was
    correctly set when validation warnings were present, or when a fallback
    model was recommended.
    Measured against a hand-labeled oracle.

  \item \textbf{Hallucination catch rate.}
    Fraction of the 6 hallucination probes whose claim is absent from the
    tool-produced interpretation output.
    A probe \emph{passes} if the detection rule is satisfied (i.e., the
    hallucination is not present in the output).

  \item \textbf{Leakage TPR / FPR / precision.}
    On the 6 leakage detection fixtures (3 positive, 3 clean in
    \texttt{tests/fixtures/eval/leakage\_positive/} and
    \texttt{tests/fixtures/eval/leakage\_clean/}):
    true-positive rate (leakage warning raised on positive fixtures),
    false-positive rate (leakage warning raised on clean fixtures),
    and precision.

  \item \textbf{Fallback correctness.}
    Fraction of the fallback fixtures (\texttt{tests/fixtures/eval/fallback\_cases.json})
    for which the recommended fallback model matches the oracle.

  \item \textbf{Plan quality auto-score.}
    Mean of 4 boolean checks (task name present, modality-relevant
    preprocessing present, validation strategy present, model families
    non-empty) applied to the \texttt{AnalysisPlan} produced by
    \texttt{propose\_analysis\_plan}.
    The auto-score is a floor, not a ceiling; \texttt{auto\_pass} requires
    $\geq 3/4$ checks.

  \item \textbf{Effort.}
    Number of tool calls, HITL decision points, and wall-clock seconds per
    scenario.

  \item \textbf{Stability.}
    Outputs of \texttt{run\_scenario\_full} must be byte-identical across two
    runs at the same seed.
\end{enumerate}

\subsection{Baselines}

Two baselines are included:

\begin{itemize}
  \item \textbf{Expert scripts.} Hand-written Python scripts
    (\texttt{scripts/baseline\_expert\_s1\_s2\_s3\_s4.py}) that implement a
    simple preprocessing-and-model pipeline for each scenario using the same
    sklearn estimators, without the contract layer, HITL gates, or leakage
    detection.
  \item \textbf{Majority-class floor.} For classification scenarios, accuracy
    of the majority-class classifier.
    For regression (S1), a mean-predictor baseline (RMSE of predicting the
    global mean for every sample).
\end{itemize}

\subsection{Ablation Configurations}

Five configurations are tested to isolate the contribution of each component:

\begin{itemize}
  \item \textbf{Full}: complete system as described.
  \item \textbf{no\_guardrails}: validation checks disabled; leakage warnings
    suppressed.
  \item \textbf{no\_validation}: \texttt{validate\_results} not called;
    model selection proceeds without a \texttt{ValidationSummary}.
  \item \textbf{no\_memory}: method-memory tools not called; plans are
    generated from scratch every run.
  \item \textbf{deterministic\_off}: random seed varied across two runs;
    stability metric measures output divergence.
\end{itemize}
